% Generated draft: requires mathematical and source review.
\documentclass[11pt]{article}
\usepackage[margin=25mm]{geometry}
\usepackage{fontspec}
\setmainfont{Latin Modern Roman}
\usepackage{amsmath,amssymb,mathtools,unicode-math}
\setmathfont{Latin Modern Math}
\usepackage{enumitem,needspace,xurl,hyperref}
\setlength{\parindent}{0pt}
\setlength{\parskip}{5pt}
\setlength{\emergencystretch}{3em}
\newcommand{\N}{\mathbb{N}}
\newcommand{\Z}{\mathbb{Z}}
\newcommand{\Q}{\mathbb{Q}}
\newcommand{\R}{\mathbb{R}}
\newcommand{\C}{\mathbb{C}}
\newcommand{\F}{\mathbb{F}}
\newcommand{\A}{\mathbb{A}}
\newcommand{\PP}{\mathbb{P}}
\newcommand{\E}{\mathbb{E}}
\newcommand{\eps}{\varepsilon}
\newcommand{\dd}{\,\mathrm{d}}
\begin{document}
\section*{UnsolvedMath ID: 30002745}
\textbf{Source problem code:} catalogue entry
\subsection{Definitions and mathematical statement}
The supplied record concerns individual students' conceptions, normative ways of mathematical reasoning that emerge in a particular classroom, and students' production-design roles. These are empirical constructs for which the record supplies no operational definitions. In particular, it does not define a production-design role or specify how coordination between contributions and roles is to be identified or measured.

The source question is retained verbatim:
\begin{quote}
What conceptions do individual students bring to their mathematical work, what normative ways of reasoning emerge in a particular classroom, and how do those contributions coordinate with the students' production-design roles?
\end{quote}

The intended setting is students' mathematical work in a particular classroom. The record does not specify the classroom, student population, mathematical activities, observation period, evidence to be collected, or criteria for identifying conceptions, normative reasoning, contributions, and coordination. It also leaves the referent of ``those contributions'' unresolved. Consequently, the supplied question does not determine a quantified mathematical proposition or a uniquely specified empirical investigation. A precise theorem, universal claim about classrooms, or numerical target cannot be extracted without additional hypotheses and definitions.
\subsection{Short English statement}
The question asks what conceptions students bring to mathematical work, which ways of reasoning become normative in a particular classroom, and how students' contributions relate to their production-design roles. Its empirical constructs and study setting require clarification before a precise mathematical or operational statement can be given. The supplied registry labels it open; that label is not independently verified.
\subsection{Sources}
\begin{itemize}
\item Local UnsolvedMath problem registry, permanent ID \(30002745\), source problem code \texttt{catalogue entry}, title ``Individual Conceptions and Classroom Mathematical Norms.'' The supplied record labels the statement exact and reports that it was retained byte-for-byte. These are supplied registry assertions, not independent source verification.
\item Supplied source citation: \emph{Mathematics in Undergraduate Study Programs: Challenges for Research and for the Dialogue between Mathematics and Didactics of Mathematics} (2015), Oberwolfach Reports, pp. 3103--3175, DOI 10.4171/owr/2014/56. Source URL and literature-source label ``original OWR report'': \url{https://doi.org/10.4171/owr/2014/56}. The supplied background describes Oberwolfach Reports as documenting mathematical workshops and their open-problem discussions. The record gives proposed year 2014 and supplies no named proposer or individual attribution for this question.
\item Additional supplied literature source: ``Rasmussen and Stephan, classroom progress study,'' \url{https://sci.sdsu.edu/crmse/msed/papers/rasmussen2.pdf}. No fuller bibliographic citation or specific supporting passage is supplied.
\item Literature assessment dated exactly 2026-08-21: the supplied curation labels the question open and describes it as a broad empirical question about individual conceptions, classroom norms, and participation. It states that existing classroom studies operationalize these constructs but do not provide a general resolution. This is supplied dated literature triage derived from the OWR curation, not an independently verified assessment. The supplied record calls for independent verification of time-sensitive status and cited claims; neither linked source has been visited for this draft.
\end{itemize}
\end{document}
